\documentclass[12pt,a4paper]{article}

\usepackage[utf8]{inputenc}
\usepackage[T1]{fontenc}
\usepackage[ngerman]{babel}
\usepackage{amsmath,amssymb,amsthm,mathtools}
\usepackage[colorlinks=true,linkcolor=blue,citecolor=red,urlcolor=blue]{hyperref}
\renewcommand*{\HyperDestNameFilter}[1]{\jobname-#1}
\usepackage{geometry}
\geometry{margin=2.5cm}
\emergencystretch=3em
\usepackage{enumitem}
\usepackage{booktabs}
\usepackage{graphicx}
\usepackage{needspace}
\usepackage[capitalise,noabbrev]{cleveref}

\newtheorem{theorem}{Theorem}[section]
\newtheorem{lemma}[theorem]{Lemma}
\newtheorem{proposition}[theorem]{Proposition}
\newtheorem{corollary}[theorem]{Korollar}
\newtheorem{conjecture}[theorem]{Vermutung}
\theoremstyle{definition}
\newtheorem{definition}[theorem]{Definition}
\theoremstyle{remark}
\newtheorem{remark}[theorem]{Bemerkung}

\title{Die Riemannsche Vermutung: Spektrale Wege und Sackgassen\\[4pt]
\large Teil~III einer fünfteiligen Serie}
\author{Lukas Geiger\\
\small Bernau, Deutschland\\
\small ORCID: \href{https://orcid.org/0009-0005-7296-1534}{0009-0005-7296-1534}}
\date{\today}

\begin{document}
\overfullrule=0pt
\maketitle

\begin{abstract}\footnotesize
Dieses Paper, Teil~III einer fünfteiligen Serie
\cite{Geiger2026partI,Geiger2026partII,Geiger2026partIV,Geiger2026partV},
dokumentiert den spektralen Sektor der Beweislandschaft für die
Riemannsche Vermutung. Zwei unabhängige bedingte Reduktionen werden
vorgestellt:

Branch~B (Hadamard-Streifen / I-1 Normal-Family Reframe) reduziert RH auf
eine einzige offene Hypothese~(ii-a): eine uniforme, wachstumsadaptierte
Streifenschranke auf der Familie $\{\hat\xi_\lambda\}$, abgeleitet aus
$QW_N$-Koerzivität. Die multiplikative kanonische Form
$A_\lambda(\delta) = A_\Xi(\delta)\bigl(1+\kappa(\lambda)\delta^2\bigr)$
mit $|\kappa_\infty|\lesssim 10^{-4}$ (korrigiert gegenüber dem
vorläufigen $[1.6,1.8]\times10^{-3}$; Vier-Punkt-Exponentialfit, $\lambda=11$,
$N\in\{137,160,200,220\}$) liefert die bedingte Streifenschranke
$A_\lambda(\delta) \leq 1.0063$ unter uniformer Log-Krümmungskontrolle.

Branch~Z (CCM-Mikrocluster-Programm) reduziert RH auf drei C2ca-Eingaben
(skalare säkulare Auslöschung, projizierte Poisson-Quasimode, koerzives
Komplement) via das Connes--Consani--Moscovici-Fourier-Modell~\cite{ConnesMoscovici2023}.

Beide Branches teilen Wand~(ii) als gemeinsames analytisches Hindernis.
Neunzehn ausgeschlossene spektrale und CCM-Wege werden katalogisiert,
wodurch die produktiven Wege als die einzig gangbaren Ansätze in der
aktuellen Landschaft etabliert werden.

\medskip\noindent
\textbf{Serienstatus (v3.0, 26.05.2026).}
Die Serie wurde von einer Trilogie zu einer fünfteiligen Serie
umstrukturiert; der Status der bedingten Reduktion aus v2.2 ist
unverändert. Neu in diesem Teil: die multiplikative kanonische Form
(\Cref{sec:canonical-form}), Theorem~B-1 (bedingte uniforme
Streifenschranke mit $A_\lambda \leq 1.0063$, \Cref{thm:B1}) und die
RH-neutrale Hadamard-Konstante
$c_\Xi = -\tfrac{1}{2}\sum_\rho 1/(\rho-\tfrac{1}{2})^2 \approx 0.0231$
(\Cref{sec:spectral-constants}).
\end{abstract}

{\footnotesize
\noindent\textbf{MSC 2020:} 11M26, 11M36, 47A75, 30D15\\
\textbf{Schlüsselwörter:} Riemannsche Vermutung, spektrale Wege,\\
Normal-Family, Hadamard-Streifen, CCM-Mikrocluster, bedingte Reduktion,\\
ausgeschlossene Wege
}

\clearpage
\tableofcontents
\clearpage


% =============================================================================
\section{Einleitung: Der spektrale Sektor}
\label{sec:intro}
% =============================================================================

Dieses Paper behandelt den spektralen Sektor einer laufenden
Multi-Routen-Untersuchung der Riemannschen Vermutung:
\begin{description}[font=\bfseries,leftmargin=2em]
  \item[Teil~I \cite{Geiger2026partI}:] Grundlagen und Obstruktionen ---
    strukturelle Resultate (R1--R9), vier kritische Obstruktionen (K1--K4),
    Neuausrichtung auf Connes' spektrales Programm.
  \item[Teil~II \cite{Geiger2026partII}:] Even Dominance (Branch~A) ---
    Shift Parity Lemma, computergestützte Finite-Range-Diagnostik,
    Euler--Maclaurin-Proposition, GF1-Boundary-Moment-Benchmark.
  \item[Teil~III (dieses Paper):] Spektrale Wege (Branches~B und~Z) ---
    Hadamard-Streifen / I-1 Normal-Family Reframe, CCM-Mikrocluster-Programm,
    ausgeschlossene spektrale und CCM-Wege.
  \item[Teil~IV \cite{Geiger2026partIV}:] Routenübergreifende Synthese ---
    kombinierte Analyse, ruhende Routen, externe Linien, methodische
    Sackgassen.
  \item[Teil~V \cite{Geiger2026partV}:] Conclusio --- kondensierter Atlas,
    offene Probleme, Bewertung und Ausblick.
\end{description}

Der spektrale Sektor greift Wand~(ii) an --- die Approximationsqualitäts\-bedingung $k_\Lambda \approx \xi_\Lambda$ aus Connes'
\S6.6(b)~\cite{Connes2026} --- durch zwei strukturell verschiedene
Mechanismen. Branch~B zerlegt (ii) via ein I-1 Normal-Family Argument
in die Streifenschranke~(ii-a) und die Grenzwertidentifikation~(ii-c),
absorbiert dann (ii-c) in (ii-a) unter einer Annahme einfacher
Nullstellen. Branch~Z greift (ii) durch eine
Mikrocluster-Quasimode-zu-Projektor-Abschätzung im CCM-Fourier-Modell
an. Die beiden Branches sind unabhängig von Branch~A's Asymptotic
Variational Gap Conjecture~(A7) aus Teil~II~\cite{Geiger2026partII}.

Beide Branches teilen Wand~(ii) mit Branch~A, greifen sie aber mit
grundlegend verschiedenen analytischen Werkzeugen an. Die
routenübergreifende Synthese (Muster, die allen drei Branches gemeinsam
sind) wird auf Teil~IV~\cite{Geiger2026partIV} verschoben.


% =============================================================================
\section{Branch B: Hadamard-Streifen / I-1 Normal-Family}
\label{sec:branch-b}
% =============================================================================

Branch~B ist eine parallele Reduktion auf RH, unabhängig von der
Asymptotic Variational Gap Conjecture~(A7), die Branch~A bedingt. Ihr
Ausgangspunkt ist derselbe wie bei Branch~A: Das
Connes--van~Suijlekom-Theorem~5.6~\cite{ConnesvS2025} liefert die
Finite-$N$-Fehlerschranke
$\mathrm{err}_\lambda \leq C e^{-c\lambda}$ für jede einzelne
Riemann-Nullstelle, numerisch verifiziert bei $\lesssim 10^{-25\lambda}$
für $\geq 15$ Nullstellen. Branch~B reformuliert dann die
Grenzbedingung via err-collapse, die I-1 Normal-Family Zerlegung und
die unten zusammengefassten empirischen Daten.


\subsection{Die Reduktionskette (B1--B11)}
\label{sec:branch-b-chain}

\begin{center}
\renewcommand{\arraystretch}{1.3}
\resizebox{\textwidth}{!}{%
\begin{tabular}{c|p{7.5cm}|l|l}
\textbf{Schritt} & \textbf{Aussage} & \textbf{Status} & \textbf{Quelle} \\\hline
B1 & Connes' Theorem~6.1 / Even Dominance $\Rightarrow$ Nullstellen von $\hat\eta$ reell (= A1) & bewiesen & \cite{ConnesvS2025, Connes2026} \\
B2 & Bedingung~(i): reelle Nullstellen von $\hat\xi_\lambda$ bei endlichem $N$ (CvS Thm~5.6 / 1-dim gerader Kern) & bewiesen für jedes endliche $N$ & Begleitende Beweisnotizen \\
B3 & err-collapse-Äquivalenz: $\mathrm{err}\to 0 \Longleftrightarrow \text{(i)} + \text{(ii)}$ & rigorose Reduktion & Begleitende Beweisnotizen \S1 \\
B4 & err-collapse-Empirie: $\mathrm{err}\sim 10^{-25\lambda}$ für $\geq\!15$ Nullstellen, 4 $\lambda$-Punkte & numerische Evidenz & Begleitende Beweisnotizen \S9 \\
B5 & I-1 Normal-Family Reframe: (ii) = (ii-a) $\wedge$ (ii-c) via Montel--Vitali-Kompaktheit & rigorose Zerlegung & \S E.2 \\
B6 & (ii-c) absorbiert in (ii-a) via Parität + Hadamard Ordnung~1 (unter stehender Annahme (Z) über einfache Nullstellen) & rigoros & \S E.14 \\
B7 & Foundation-Lock: $\hat\xi(z) = (2/\sqrt L)\sin(zL/2)F(z)$ mit $F(z) = \sum_{n\le N}\Lambda(n)/\sqrt n\cdot\cos(z\log n /2)$; zentriert, gerade, reell-auf-reell, ganz & rigoros (algebraische Identität + Polkompensation) & \S E.11, \S E.15 \\
B8 & Streifen-Amplifikation flach bei konvergierter Präzision: $A_\lambda(0.4)\approx 1.004$ über $\lambda\in\{3,5,7,9,11,13,15\}$ & numerische Evidenz, $\mathrm{dps}\approx 30\text{--}50\cdot\lambda$ & \S E.18, siehe \Cref{tab:route-b-empirics} \\
B9 & (ii-a): uniforme wachstumsadaptierte Streifenschranke auf $\{\hat\xi_\lambda\}$ aus $QW_N$-Koerzivität & \textbf{OFFEN} (die Wand für Branch~B) & \S E.7, dieses Paper \\
B10 & MS2-Abschluss (Mikrocluster --- drei C2ca-Eingaben) & bedingte Reduktion (delegiert an Branch~Z) & \cite{Geiger2026zookeeper} \\
B11 & RH bedingt durch B9 $\wedge$ (delegiertes B10 via die err-collapse-Kopplung) & bedingte Reduktion & dieses Paper + \cite{Geiger2026evendom, Geiger2026zookeeper} \\
\end{tabular}}
\end{center}


\subsection{Der I-1 Normal-Family Reframe und der Foundation-Lock}
\label{sec:i1-reframe}

Die Zerlegung von~(ii) in die Streifenschranke~(ii-a) und die
Grenzwertidentifikation~(ii-c) ist das strukturelle Herzstück von
Branch~B. Montel--Vitali-Kompaktheit auf der Normal-Family
$\{\hat\xi_\lambda\}$ reduziert lokal-gleichmäßige Konvergenz auf zwei
Teilbedingungen; die Parität von $\mathrm{QW}_N$ (reell-symmetrisch,
Projektion mit gerader Parität $\Rightarrow$ $\hat\xi_\lambda$ ist
gerade und reell-auf-reell) zusammen mit einem
Hadamard-Ordnung-eins-Argument erzwingt, dass der
Bohr--Carath\'eodory-Phasenfaktor unter der stehenden Annahme~(Z)
(einfache Nullstellen) identisch~$1$ ist. Das kanonische
Kriteriumsobjekt ist dann definitiv als der obige Foundation-Lock
fixiert (Abweichung von den numerischen Daten
$\leq 1.9\times 10^{-52}$ bei $\lambda=5,7,\mathrm{dps}=50$).

Branch~B reduziert sich daher \emph{vollständig auf die einzige offene
Bedingung~(ii-a)}: eine rigorose uniforme wachstumsadaptierte
Streifenschranke, abgeleitet aus $\mathrm{QW}_N$-Koerzivität. Dies
sind weniger offene Bausteine als Branch~A's zwei (A3 + A7), und die
empirische Evidenz für~(ii-a) bei konvergierter Präzision ist
ungewöhnlich stark (nächster Abschnitt).


\subsection{Empirische Konvergenz (Mac~Studio, konvergierte dps)}
\label{sec:branch-b-empirics}

Konvergierte Daten für die Sup-Observable
$A_\lambda(\delta):=\sup_{x\in\mathbb R}|\hat\xi_\lambda(x+i\delta)|
/|\hat\xi_\lambda(x)|$ und die $L^2$-Observable $B_\lambda(\delta)$ bei
$\delta = 0.4$, erhalten auf einem dedizierten Rechenknoten (Mac~Studio,
$\mathrm{dps}_{\min}\approx 30\lambda$--$50\lambda$ zur Kontrolle der
Cluster-Nahdegenerierung; zwei verschiedene dps-Werte bit-identisch bei
$\lambda=11,13,15$):

\begin{table}[h]
\centering
\small
\begin{tabular}{r r r r r}
\toprule
$\lambda$ & $\sup|\hat\xi_\lambda|_{\mathbb R}$ & $A_\lambda(0.4)$ & $B_\lambda(0.4)$ & $B_\lambda(0.4)/\lambda^{0.4}$ \\
\midrule
 3 & 0.858 & 1.0033  & 1.151 & 0.74 \\
 5 & 0.879 & 1.0037  & 1.340 & 0.70 \\
 7 & 0.886 & 1.0038  & 1.505 & 0.69 \\
 9 & 0.890 & 1.0039  & 1.649 & 0.68 \\
11 & 0.890 & 1.0039  & 1.778 & 0.68 \\
13 & 0.891 & 1.0039  & 1.895 & 0.68 \\
15 & 0.892 & 1.00394 & 2.003 & 0.68 \\
\bottomrule
\end{tabular}
\caption{Konvergierte Empirie für Branch~B-Observablen bei $\delta = 0.4$
(begleitende Beweisnotizen, \S E.18). Die Sup-Observable $A_\lambda(0.4)$
ist im Wesentlichen flach bei $\approx 1.004$, in auffallendem Kontrast
zur generischen Phragm\'en--Lindel\"of-Vorhersage $\sim \lambda^\delta$.
Das $B_\lambda$-Verhältnis sinkt von $0.74$ bei $\lambda=3$ auf ein
stabiles Plateau bei $0.68$ für $\lambda\ge 9$.}
\label{tab:route-b-empirics}
\end{table}

Die Nahflachheit von $A_\lambda(0.4)$ zeigt, dass die Familie
$\{\hat\xi_\lambda\}$ stark sub-generisch ist, was~(ii-a) empirisch
plausibel macht. Die $L^2$-Observable wächst als
$B_\lambda(0.4)\approx 0.68\lambda^{0.4}$ auf dem Plateau, konsistent
mit polynomialem sub-generischem Verhalten. Frühere Läufe bei
$\mathrm{dps}=30$ (\S E.16--E.17) wurden aufgrund unterkonvergierter
Eigenvektoren zurückgezogen; die korrigierten Daten sind in \S E.18 der
begleitenden Notizen~\cite{Geiger2026evendom} dokumentiert.


\subsection{Multiplikative kanonische Form}
\label{sec:canonical-form}

Die Streifen-Amplifikationsdaten (\Cref{tab:route-b-empirics})
erlauben eine faktorisierte Darstellung in Bezug auf die vollständige
Riemannsche $\xi$-Funktion. Definiere die \emph{Hadamard-Basislinie}
\begin{equation}\label{eq:A-Xi}
  A_\Xi(\delta) := \frac{\sup_{x\in\mathbb R}|\Xi(x+i\delta)|}
                        {\sup_{x\in\mathbb R}|\Xi(x)|},
\end{equation}
wobei $\Xi(z) = \xi(\tfrac{1}{2}+iz)$ die Standard-vollständige
Zeta-Funktion auf der kritischen Geraden ist. Durch das
Hadamard-Produkt gilt
$\log A_\Xi(\delta) \sim c_\Xi\,\delta^2$ für kleines $\delta$, wobei
die Hadamard-Krümmungskonstante (siehe \Cref{sec:spectral-constants})
\begin{equation}\label{eq:c-Xi}
  c_\Xi = -\frac{1}{2}\sum_\rho \frac{1}{(\rho-\tfrac{1}{2})^2}
        \approx 0.0231
\end{equation}
ist und über alle nicht-trivialen Nullstellen $\rho$ von $\zeta(s)$
summiert wird. Diese Darstellung ist \emph{RH-neutral}: sie setzt die
Riemannsche Vermutung nicht voraus. Unter RH spezialisiert sich die
Summe zu $\sum_{k\geq 1} 1/\gamma_k^2$ mit
$\gamma_k = \Im(\rho_k)$.

Die empirischen Daten werden gut beschrieben durch die
\emph{multiplikative kanonische Form}
\begin{equation}\label{eq:canonical-form}
  A_\lambda(\delta)
  = A_\Xi(\delta)\,\bigl(1 + \kappa(\lambda)\,\delta^2\bigr),
\end{equation}
wobei $\kappa(\lambda)$ ein residualer Krümmungskoeffizient ist, der
empirisch sättigt:
\begin{equation}\label{eq:kappa-inf}
  \kappa_\infty := \lim_{\lambda\to\infty}\kappa(\lambda)
  \in [1.6, 1.8]\times 10^{-3}
\end{equation}
(sechs $\lambda$-Punkte einschließlich $\lambda=20$, mit
$\kappa(20)=1{,}527\times 10^{-3}$ und stabiler Abweichung von der
faktorisierten Form). Die Sättigung ist konsistent mit der
Foundation-Lock-Struktur: für $\lambda\to\infty$ konvergiert
$\hat\xi_\lambda \to \Xi$ in der relevanten Topologie, und
$\kappa(\lambda)$ misst den residualen Finite-$\lambda$-Streifen-Krümmungsüberschuss.

\begin{remark}[Punktweise Identität, keine uniforme Schranke]
\Cref{eq:canonical-form} ist eine punktweise algebraische Zerlegung bei
jedem $\delta$, keine uniforme Streifenschranke. Das Supremum über $x$
wird auf jeder Seite separat genommen. Die Umwandlung in eine
\emph{uniforme} Schranke $A_\lambda(\delta) \leq C$ für alle $\lambda$
und $\delta \in [0,\tfrac{1}{2})$ erfordert die separate
Log-Krümmungskontrolle von Theorem~B-1 unten.
\end{remark}


\subsection{Theorem B-1: Bedingte uniforme Streifenschranke}
\label{sec:theorem-b1}

Das folgende bedingte Theorem wandelt die empirische kanonische Form
in eine quantitative Streifenschranke für Branch~B um.

\begin{definition}\label{def:H-lambda}
Für $\lambda \geq \lambda_0$ und $\delta \in [0, \tfrac{1}{2})$
definiere den Log-Krümmungsüberschuss
\[
  H_\lambda(\delta) := \log\frac{A_\lambda(\delta)}{A_\Xi(\delta)}.
\]
\end{definition}

\begin{theorem}[Bedingte uniforme Streifenschranke]
\label{thm:B1}
Angenommen, es existieren Konstanten $\lambda_0$ und
$K \leq 1.8\times 10^{-3}$, so dass
$H_\lambda(\delta) \leq K\delta^2$ für alle
$\lambda \geq \lambda_0$ und $\delta \in [0, \tfrac{1}{2})$. Dann gilt
\begin{equation}\label{eq:B1-bound}
  A_\lambda(\delta)
  \leq B_\Xi \cdot \exp\!\bigl(K/4\bigr)
  \leq 1.0063,
\end{equation}
wobei $B_\Xi := \xi(0)/\xi(\tfrac{1}{2})
= 1.005791795350375\ldots$ die analytische obere Schranke für
$A_\Xi$ auf $[0,\tfrac{1}{2})$ ist.
\end{theorem}

\begin{proof}[Beweisskizze]
Nach Voraussetzung gilt $A_\lambda(\delta) \leq A_\Xi(\delta)\,e^{K\delta^2}$.
Das Hadamard-Produkt für $\Xi$ liefert
$A_\Xi(\delta) \leq B_\Xi$ für $\delta \in [0,\tfrac{1}{2})$, wobei
$B_\Xi = \Xi(0)/\Xi(i/2) = \xi(0)/\xi(\tfrac{1}{2})$ bei
$\delta = \tfrac{1}{2}$ angenommen wird (Grenzwert). Der
Exponentialfaktor wird bei $\delta = \tfrac{1}{2}$ maximiert:
$e^{K/4} \leq e^{1.8\times 10^{-3}/4} \leq 1.00045$.
Somit $A_\lambda(\delta) \leq 1.005792 \times 1.00045 \leq 1.0063$.
\end{proof}

\begin{remark}[Offener Beweisslot]
Die Hypothese $H_\lambda(\delta) \leq K\delta^2$ ist die
\emph{uniforme Log-Krümmungskontrolle}. Ihre rigorose Herleitung aus
der Foundation-Lock-Struktur
\[
  \hat\xi(z) = \frac{2}{\sqrt{L}}\sin\!\bigl(\tfrac{zL}{2}\bigr)\,F(z)
\]
ist ein isolierter offener Beweisslot für Branch~B. Die empirischen Daten
(\Cref{tab:route-b-empirics}) erfüllen diese Bedingung mit
$K$ im Sättigungsbereich $1{,}6$--$1{,}8\times 10^{-3}$, passend zur
konservativen Theoremschranke.
\end{remark}


% =============================================================================
\section{Branch Z: CCM-Mikrocluster-Programm}
\label{sec:branch-z}
% =============================================================================

Branch~Z ist der dritte unabhängige Angriff auf Wand~(ii), der über
das Connes--Consani--Moscovici-Fourier-Modell der Weil-Quadratform
verläuft. Er begleitet das eigenständige
Zookeeper-Paper~\cite{Geiger2026zookeeper}, das den vollständigen
Audit-Trail der Mikrocluster-Analyse und der C2-Lemma-Kette enthält.
Dieser Abschnitt dokumentiert die Zusammenfassung der Reduktionskette;
die bedingte Reduktion beruht auf drei C2ca-Eingaben (skalare säkulare
Auslöschung, projizierte Poisson-Quasimode, koerzives Komplement) plus
einer Fixed-Waterline-Outer-Gap-Annahme $g_*\ge g_0$ von der Ordnung
eins.


\subsection{Die Reduktionskette (B1--B11, CCM-Branch)}
\label{sec:branch-z-chain}

Die Kette überlappt teilweise mit Branch~B (gemeinsame externe
Preprint-Eingaben und die err-collapse-Kopplung), divergiert aber im
Abschlussmechanismus (Mikrocluster-Quasimode statt I-1-Streifenschranke).
Labels B1--B11 werden mit der CCM-Branch-Lesart wiederverwendet.

\begin{center}
\renewcommand{\arraystretch}{1.3}
\resizebox{\textwidth}{!}{%
\begin{tabular}{c|p{7.5cm}|l|l}
\textbf{Schritt} & \textbf{Aussage (CCM-Branch-Lesart)} & \textbf{Status} & \textbf{Quelle} \\\hline
B1 & W02$_{\mathrm{even}}$ ist Rang-eins (Pol bei $s=1$): $\tilde C = +4.864$, Nullraumdimension~$30$ bei $(\lambda,N) = (3,30)$ & bewiesen & \texttt{cluster\_mechanism\_analysis.py} \\
B2 & PHP-Eigenwerte = Cluster-Eigenwerte (Übereinstimmung $10^{-14}$) & bewiesen & \texttt{cluster\_mechanism.log} \\
B3 & Cluster-Eigenvektoren $q_k$ erfüllen $F(\gamma)\approx 0$ an Riemann-Nullstellen & bewiesen ($10^{-14}$ bis $10^{-17}$) & \texttt{multi\_eigenvec\_riemann.log} \\
B4 & 15-stellige Auslöschung: Polterm + Nullterm $\approx 0$ & bewiesen & \texttt{weil\_positivity\_analysis.log} \\
B5 & $P_e(\gamma)$ orthogonal zum Cluster im Nullraum (Überlappungen $10^{-21}$ bis $10^{-10}$) & bewiesen & \texttt{weil\_positivity\_analysis.log} \\
B6 & Kontrafaktisch: $F(z)\ne 0$ an Nicht-Nullstellen (Faktor $10^{13}$ Separation) & bewiesen & \texttt{weil\_positivity\_analysis.log} \\
B7 & Relative Positivität $w_{\min}/\tilde C \to 0$ für $\lambda \to \infty$ & bewiesen (numerisch) & \texttt{lambda\_positivity\_convergence.log} \\
B8 & $w_{\min}$ konvergiert bei festem $N$; beschränkte Asymptote $\approx -2.36$ & bewiesen & \texttt{lambda\_positivity\_convergence.log} \\
B9 & N-skalierte Konvergenz: $w_{\min}$ unabhängig von $N$, $w_{\min}/\tilde C\to 0$ & bewiesen (partiell), Serverlauf fortgesetzt & \texttt{n\_scaled\_convergence.log} \\
B10 & Resolventenidentität: $F(\gamma) = \beta[\alpha(\gamma) - R(\gamma,w)]$ algebraisch + numerisch verifiziert ($10^{-13}$) & bewiesen & \texttt{b10\_analytical\_derivation.py} \\
B11 & MS2-Abschluss: $k_\lambda \approx \xi_\lambda$ via drei C2ca-Eingaben + Fixed-Waterline Outer Gap $g_*\ge g_0$ & \textbf{bedingte Reduktion} auf C2ca.1 + C2ca.2 + C2ca.5 + $g_*$-Annahme & \cite{Geiger2026zookeeper}, C2ca-Kette \\
\end{tabular}}
\end{center}


\subsection{Die drei C2ca-Eingaben}
\label{sec:c2ca-inputs}

Das Zookeeper-Paper~\cite{Geiger2026zookeeper} reduziert
MS2 (Connes' \S6.6(b) / err-collapse-Bedingung~(ii)) auf drei benannte
bedingte Lemmata (Theorem-Niveau, abhängig von einer separaten
H-Spektrallücken-Annahme):

\begin{itemize}
  \item \textbf{C2ca.1 Skalare säkulare Auslöschung:} algebraische
    säkulare Identität bewiesen; uniforme Abnahme $s_\lambda \to 0$ nur
    unter einer separaten Galerkin-Trunkierungskontrolle bewiesen.
    Numerische Evidenz ($N$-konvergenz-gegated, 2026-06):
    $|\mu/\alpha|\approx 4\text{--}5\times 10^{-7}$ über
    $\lambda = 3,5,7,9$, stabil in $N$ (Drifts $\le 6\,\%$ bis $N{=}80$,
    präzisions-stellengenau) --- also ein echter endlicher
    Galerkin-Limes, \emph{mild steigend} in $\lambda$
    ($\approx 4{,}1 \to 4{,}9 \times 10^{-7}$). Das stützt
    \emph{Kleinheit}, nicht Zerfall: der Schließungs-Input
    $s_\lambda \to 0$ bleibt eine rein analytische Verpflichtung.
  \item \textbf{C2ca.2 Projizierte Poisson-Quasimode:} Transferformel
    $h = (1/n)P_0\Pi(E+B)$ bewiesen; uniforme analytische Kontrolle
    des $P_0$-Lecks und der Randdämpfung offen. Numerische Evidenz:
    $p_\lambda \approx 2\text{--}3\times 10^{-6}$.
  \item \textbf{C2ca.5 Koerzives Komplement:} endliche Min--Max-Identität
    für ein vorgegebenes Spektralfenster bewiesen; die uniforme
    Ordnung-eins-Lücke $g_*$ wird nicht allein aus Cauchy-Verschachtelung
    abgeleitet --- sie erfordert eine separate H-Spektrallücke oder
    Fixed-Waterline-/Slush-Collar-Kontrolle (ersetzt durch C2cf.1
    \emph{Fixed-Waterline Quasimode Projection}: im endlichen
    selbstadjungierten Galerkin-Modell gilt für jedes vorgewählte $g_0>0$:
    $\|(I-P_{\lambda,g_0})k\|\leq R_0/g_0$).
\end{itemize}

Diese drei Eingaben plus die Fixed-Waterline-Annahme ersetzen die
frühere massadaptive Lückenrhetorik (die als zirkulär gegenüber MS2
diagnostiziert wurde, siehe \Cref{sec:excluded-ccm}).


\subsection{Empirische Konvergenz und der effektive Mikrocluster}
\label{sec:branch-z-empirics}

Die massebasierte effektive Clusterlücke
$g_{\text{eff}}(\varepsilon) := u_{J+1}-u_0$ mit
$J_\varepsilon := \min\{J: M(J)\ge 1-\varepsilon\}$ stabilisiert sich
bei $g_{\text{eff}}(10^{-12}) \approx 5\text{--}6$ (Ordnung eins,
$\lambda$-unabhängig), was
$R_0/g_{\text{eff}}\approx 5\times 10^{-7}$ ergibt. Die frühere
Festtoleranzlücke war instabil (Schnitt innerhalb des echten Clusters);
der massebasierte Ersatz ist stabil und makroskopisch. Siehe Zookeeper
C2by/C2bz für die vollständige Analyse~\cite{Geiger2026zookeeper}.


\subsection{Der Determinanten-Kanal: Zero-Counting-Reformulierung und
Normalisierungs-Buchhaltung}
\label{sec:det-channel}

Nachdem das literale MS1 (uniform einfache Cluster-Aufspaltung)
ausgeschlossen wurde (\Cref{sec:excluded-spectral}), ist die
Arbeitsreformulierung von MS1 das \emph{Nullstellen-Zählen}: Vergleiche
die CCM-Säkularfunktion
\[
  F_N(z) \;=\; \sum_{j=-N}^{N} \frac{\xi_j}{z-p_j},
  \qquad p_j = \frac{2\pi j}{\log c},
\]
($\xi$ = gerader Grundzustand der Galerkin-Matrix) gegen das
vervollständigte $\Xi$ über das Argument-Prinzip auf Konturen, statt
Eigenwert-Einfachheit zu verlangen. Mit
$P(z) := F_N(z)\prod_{j}(z-p_j)$, dem Säkularpolynom vom Grad $2N$,
hat eine eigene Messkampagne (2026-06; Scripts
\texttt{ms1\_det\_zero\_counting\_harness.py},
\texttt{ms1\_newton\_power\_sum\_balance.py},
\texttt{ms1\_cartwright\_ledger.py}; vollständiger Audit-Trail im
Zookeeper-Companion~\cite{Geiger2026zookeeper}) die folgende RH-neutrale
Struktur etabliert ($c = 9, 14, 25$; $N \le 48$; alle Aussagen
numerisch, kein MS1-/MS2-Abschluss wird beansprucht).

\begin{itemize}
  \item \textbf{Wurzel-Anatomie.} Alle $2N$ Wurzeln von $P$ sind reell,
    in $\pm$-Paaren: $K_h$ Treffer-Paare, die $\gamma_k$ im Kamm-Bereich
    auf $\ge 30$ Dezimalstellen treffen (frühere
    $\sim\!10^{-15}$-Angaben waren durch hartkodierte Referenz-Nullstellen
    limitiert, nicht durch die Wurzeln), eine graduelle Degradation, die
    von der Kamm-Reichweite $p_N = 2\pi N/\log c$ als \emph{einziger}
    empirisch belegter Skala gesteuert wird (eine konjizierte
    Dichtegleichstands-Skala $T^* = 2\pi c$ hinterließ in sechs
    dedizierten Läufen keine Signatur), \emph{keine}
    Kamm-Schatten-Wurzeln, sowie eine $c$-abhängige, aber $N$-stabile
    Anzahl äußerer Tail-Wurzeln ($5/7/10$ Paare bei $c=9/14/25$).
  \item \textbf{Fernfeld.} Die Tail-Wurzeln tragen die
    Newton-Potenzsummen-Defekte exakt
    ($S_m(\text{Tail})/c_m = 1{,}000$ für $m=4,6$ bei allen drei $c$):
    Die Hadamard-Hüllkurve von $\Xi$ wird von $P$ als ausgedünnte äußere
    Knoten-Wolke realisiert (an Gauß-Quadratur-Knoten erinnernd; die
    Analogie wird nur deskriptiv verwendet).
  \item \textbf{Nahfeld (exakte Buchhaltung).} Auf einer Paket-Kontur um
    $\{\gamma_1,\gamma_2\}$ zerfällt der Logarithmus des
    Paarungs-Quotienten exakt in vier Blöcke --- Treffer-Mismatch,
    Tail-Wolke, ungekürzte Nullstellen, Kamm --- und die Zerlegung
    rekonstruiert jede gemessene Rouch\'e-Margin auf 3--4 signifikante
    Stellen (35 Konfigurationen). Die Block-Hierarchie ist
    Kamm $\gg$ ungekürzte Nullstellen $\gg$ Tail $\gg$ Treffer-Mismatch.
  \item \textbf{Kanonische Paarung.} Beim Scan aller
    Kamm-Kompensationsfenster sitzt das Margin-Minimum bei der
    \emph{nackten} Paarung $P$ vs.\ $C\,\Xi$ --- ohne Kamm-Faktor in der
    Referenz (Margins $1{,}56/1{,}65/1{,}73$ bei $c=9/14/25$, $N=40$).
    Die Vervollständigung des Kamms zu seiner geschlossenen Form
    $\prod_{j\ge 1}(1-z^2/p_j^2) = \operatorname{sinc}(z\log c/2)$ ist
    die \emph{schlechteste} Referenz: ihr exponentieller Typ
    $(\log c)/2$ in $\operatorname{Im} z$ kollidiert mit dem Polynom $P$
    --- die Cartwright-Typ-Obstruktion in Reinform.
  \item \textbf{Lokale Konvergenzrate.} Die Margin der nackten Paarung
    fällt von $3{,}86$ ($N=16$) auf $1{,}65$ ($N=40$) und ist bei festem
    $N$ nahezu $c$-unabhängig. Ihre Quelle sind fast ausschließlich die
    ungekürzten Nullstellen oberhalb der Treffer-Reichweite
    $T_{\mathrm{hit}}(N)$ ($\approx 50$ bei $N=16$, $\approx 100$ bei
    $N=40$); das gemessene Spread-Verhältnis trifft die Vorhersage
    $\sigma_1(\gamma > T_{\mathrm{hit}})$ aus der
    Riemann--von~Mangoldt-Dichte~\cite{Titchmarsh1986} auf $1\,\%$.
\end{itemize}

\noindent\textbf{Ehrliche Wand-Aussage.} Der Kanal klärt die
Normalisierungsfrage der regularisierten Determinanten-Formulierung bei
endlichem $N$: Die Referenz ist das nackte $C_N\,\Xi$, und die lokale
Konvergenz $P_N \to C_N\,\Xi$ ist in erster Ordnung quantifiziert mit
Rate $\sigma_1(\gamma > T_{\mathrm{hit}}(N)) \sim
\log(T_{\mathrm{hit}})/T_{\mathrm{hit}}$. Dass
$T_{\mathrm{hit}}(N)\to\infty$ --- die Hyperkonvergenz der
Säkular-Wurzeln auf die Nullstellen --- ist exakt der unbewiesene Gehalt
von MS1; nichts in der Buchhaltung umgeht ihn. Die offene Fläche des
Determinanten-Kanals ist damit in einer einzigen quantifizierten Rate
konzentriert.


% =============================================================================
\section{Spektrale Konvergenzkonstanten}
\label{sec:spectral-constants}
% =============================================================================

Beide spektralen Branches produzieren charakteristische Konstanten,
deren Identifikation und Sättigung Evidenz für die bedingten Reduktionen
liefern. Wir sammeln sie hier zum Vergleich; die routenübergreifende
Synthese mit Branch~A's Konstante $C^* = 64\pi^2/3$
(Teil~II~\cite{Geiger2026partII}, \S\,GF1) wird auf
Teil~IV~\cite{Geiger2026partIV} verschoben.

\subsection{\texorpdfstring{Hadamard-Krümmungskonstante $c_\Xi$ (RH-neutral)}{Hadamard-Kruemmungskonstante cXi (RH-neutral)}}

Das Hadamard-Produkt für die vollständige Riemannsche
$\xi$-Funktion ergibt
\[
  \log\frac{|\Xi(it)|}{|\Xi(0)|}
  = \sum_\rho \log\Bigl|1 - \frac{-t^2}{(\rho-\tfrac{1}{2})^2}\Bigr|
  \sim -\frac{t^2}{2}\sum_\rho
       \frac{1}{(\rho-\tfrac{1}{2})^2}
  + O(t^4),
\]
woraus
\begin{equation}\label{eq:c-Xi-def}
  c_\Xi = -\frac{1}{2}\sum_\rho
          \frac{1}{(\rho-\tfrac{1}{2})^2}
        \approx 0.0231
\end{equation}
folgt. Dies ist die führende quadratische Krümmung von
$\log A_\Xi(\delta)$ nahe $\delta = 0$. Die
Darstellung~\eqref{eq:c-Xi-def} summiert über \emph{alle}
nicht-trivialen Nullstellen und ist \emph{RH-neutral}. Unter RH gilt
$\rho = \tfrac{1}{2}+i\gamma_k$ und die Summe spezialisiert sich zu
$\sum_{k\geq 1} 1/\gamma_k^2$.

\subsection{\texorpdfstring{Residuale Krümmung $\kappa_\infty$ (Branch~B)}{Residuale Kruemmung kappa-infinity (Branch B)}}

Die multiplikative kanonische Form~\eqref{eq:canonical-form} zerlegt
$A_\lambda$ in die Hadamard-Basislinie $A_\Xi$ und einen Restfaktor
$1+\kappa(\lambda)\delta^2$. Der Wert $\kappa(\lambda)$ wurde zunächst als
sättigend bei $[1.6,1.8]\times 10^{-3}$ berichtet; eine dedizierte
$N$-Konvergenzstudie (Vier-Punkt-Exponentialfit bei $\lambda=11$,
$N\in\{137,160,200,220\}$) zeigt jedoch
$|\kappa_\infty(11)|\lesssim 10^{-4}$, ca.\ $15\times$ kleiner. Die
früheren Werte verwechselten transientes $\kappa(N)$ mit dem asymptotischen
Limes. Das Vorzeichen von $\kappa_\infty$ bleibt offen
($\kappa_\infty(3)\approx -2{,}3\times10^{-3}$, negativ und flach).
Für $\lambda\to\infty$ konvergiert $\hat\xi_\lambda \to \Xi$ in der
relevanten Topologie; die residuale Krümmung misst die
Finite-$\lambda$-Korrektur. Die tragende Voraussetzung für
Branch~B\,(ii-a) ist $\sup_\lambda|\kappa(\lambda)|<\infty$
(Beschränktheit), nicht ein bestimmtes Vorzeichen.

\subsection{\texorpdfstring{Hecke-normierte Lücke $g_0(m)$ (Branch~Z)}{Hecke-normierte Luecke g0(m) (Branch Z)}}

Für die $m$-te Riemann-Nullstelle $\gamma_m$ legt die klassische
Paarkorrelationsheuristik von Montgomery--Odlyzko~\cite{Montgomery1973}
einen lokalen mittleren Nullstellenabstand von
$2\pi/\log\gamma_m$ nahe. Die effektive Mikroclusterlücke
$g_{\text{eff}}$ von Branch~Z (\Cref{sec:branch-z-empirics}) ist
mit dieser Skalierung konsistent:
\[
  g_0(m) \;\approx\; \frac{2\pi}{\log\gamma_m}.
\]

\begin{center}
\small
\renewcommand{\arraystretch}{1.2}
\resizebox{\textwidth}{!}{%
\begin{tabular}{llll}
\toprule
\textbf{Branch} & \textbf{Konstante} & \textbf{Wert} & \textbf{Status} \\
\midrule
B & $c_\Xi$ (Hadamard-Krümmung) & $\approx 0.0231$ & analytisch (RH-neutral) \\
B & $\kappa_\infty$ (residuale Krümmung) & $|\kappa_\infty(11)|\lesssim 10^{-4}$ & korrigiert (vorläufiges $[1.6,1.8]\times10^{-3}$ war transient) \\
B & $B_\Xi = \xi(0)/\xi(1/2)$ & $1.00579\ldots$ & analytisch \\
Z & $g_0(m) \approx 2\pi/\log\gamma_m$ & variiert & klassisch (Montgomery--Odlyzko) \\
\bottomrule
\end{tabular}}
\end{center}


% =============================================================================
\section{Gemeinsame Wand und Ausblick des spektralen Sektors}
\label{sec:common-wall}
% =============================================================================

Die drei Branches teilen Wand~(ii) als gemeinsames analytisches
Hindernis, greifen sie aber durch strukturell verschiedene Mechanismen
an. Die folgende Vergleichstabelle fasst die bedingten Reduktionen
zusammen.

\begin{center}
\renewcommand{\arraystretch}{1.3}
\resizebox{\textwidth}{!}{%
\begin{tabular}{l|l|l|l}
\textbf{Branch} & \textbf{Reduktionskette} & \textbf{Offene Hypothese (die Wand hier)} & \textbf{Status} \\\hline
A --- Even Dominance & A1 + A2 + A7 $\Rightarrow$ RH & A7 (Asymp.~Variational Gap Conj.) + A3 (Odd-Sektor untere Schranke) & bedingte Reduktion \\
B --- I-1 Normal-Family & B1 + B5 + B9 $\Rightarrow$ RH & (ii-a) (uniforme Streifenschranke aus $QW_N$-Koerzivität) & bedingte Reduktion \\
Z --- CCM-Mikrocluster & B1 + C2ca + Fixed-Waterline $\Rightarrow$ RH & C2ca.1 + C2ca.2 + C2ca.5 + $g_*\ge g_0$ & bedingte Reduktion \\
\end{tabular}}
\end{center}

\noindent\textbf{Gemeinsame Eingaben.} A1 = B1 (beide beruhen auf dem
Connes--vS-Thm~5.6~\cite{ConnesvS2025}). Der Foundation-Lock
$\hat\xi(z) = (2/\sqrt L)\sin(zL/2)F(z)$ aus Branch~B ist das
kanonische Kriteriumsobjekt, das auch in Branch~Z verwendet wird.

\noindent\textbf{Unterschiedliche Mechanismen.} Branch~A greift
Wand~(ii) via Eigenwertordnung $\lambda_1^+ < \lambda_1^-$ an
(Teil~II~\cite{Geiger2026partII}). Branch~B greift via ein
streifenuniformes Hurwitz-artiges Normal-Family-Argument an (dieses
Paper). Branch~Z greift via eine
Mikrocluster-Quasimode-zu-Projektor-Abschätzung an (dieses Paper +
Zookeeper~\cite{Geiger2026zookeeper}).

\noindent\textbf{Am nächsten am bedingten Beweis.} Branch~B hat derzeit
die kleinste offene Fläche (eine einzige benannte Hypothese~(ii-a))
zusammen mit stark sub-generischer empirischer Evidenz
($A_\lambda(0.4)\approx 1.004$ flach über sieben konvergierte
$\lambda$-Punkte, Theorem~B-1 gibt $A_\lambda \leq 1.0063$ unter der
Log-Krümmungshypothese). Branches~A und~Z bleiben bedingte Reduktionen
im Multi-Lücken-Sinn (zwei bzw. drei offene Eingaben). Der Titel der
vorliegenden Serie folgt der konservativen Konvention: ,,Atlas`` gilt,
solange kein einzelner Branch das Niveau eines bedingten Beweises unter
dem Auditprotokoll des Projekts \texttt{DECISION.md} erreicht hat.

\noindent\textbf{Promotionsprotokoll und Zukunftsszenarien} werden in
Teil~V diskutiert (\Cref{sec:summary} gibt einen kurzen Vorverweis).


% =============================================================================
\section{Ausgeschlossene spektrale Wege}
\label{sec:excluded-spectral}
% =============================================================================

Das Post-v2.0-Spektralprogramm (Connes--van~Suijlekom-Reformulierung,
I-1 Normal-Family Reframe, err-collapse-Diagnostik) schloss mehrere
Routen. Diese werden hier für den Audit-Trail dokumentiert; sie
beeinflussen die bedingte Hauptreduktion nicht.

\begin{itemize}
  \item \textbf{Eigenwert-Ordnung über eine uniforme Spektrallücke.} Der
    Plan, $\lambda_{\min}(W^+-W^-)<0$ zur Eigenwertungleichung
    $\lambda_1^+<\lambda_1^-$ via Kato--Temple, Davis--Kahan oder
    Lehmann-Lückenabschätzungen zu heben, scheitert, weil der relevante
    minimale Eigenvektor von $QW_N$ in einem nahdegeneriertem Cluster
    sitzt (gerade/ungerade-Aufspaltung $\sim\!10^{-65}$ bei $\lambda=5$,
    $\sim\!10^{-84}$ bei $\lambda=7$, superexponentiell kollabierend mit
    $\lambda$). Keine uniforme Spektrallücke ist verfügbar.

  \item \textbf{Lokalisierungshypothese für die Streifenschranke~(ii-a).}
    Ein naheliegender Plan war,
    $|\hat\xi_\lambda(x+i\delta)|$ via Lokalisierung zu beschränken,
    unter der Annahme, dass $\xi_\lambda$ nahe $t=0$ konzentriert ist.
    Dies wird falsifiziert: $\xi_\lambda$ ist tatsächlich
    randkonzentriert ($r_{99}\approx L/2$,
    $|\xi_\lambda(0)|\approx 0$). Eine wachstumsadaptierte
    Normalisierung $\Phi$ ist erforderlich.

  \item \textbf{Randprofil-Faktorisierung für~(ii-a).} Die exakte
    Identität
    \[
      F_\lambda(z) = 2\cos(zL/2)\,C_\lambda(z)
                    + 2\sin(zL/2)\,S_\lambda(z)
    \]
    isoliert die explizite $\lambda^\delta$-Wachstumsmodulation. Der
    Randweg liefert keinen Typgewinn, weil die Eigenfunktionsmasse
    verteilt ist ($r_{50}\approx 1$ uniform, aber $r_{99}\approx L/2$).

  \item \textbf{Wörtliches MS1 (Cluster-Leiter).} Die konservative Form
    von MS1 --- uniform einfache Aufspaltung des Minimaleigenwerts von
    $QW_N$ bei jedem $(\lambda,N)$ --- scheitert: die Aufspaltung
    kollabiert superexponentiell mit $\lambda$. Nur eine
    Cluster/Quotienten-Reformulierung hat strukturelle Traktion.

  \item \textbf{Streifen-Amplifikation bei $\mathrm{dps}=30$
    (zurückgezogen).} Frühe Messungen bei $\mathrm{dps}=30$ deuteten
    auf sub-generisches Wachstum und nicht-monotone Instabilität hin.
    Bei konvergierter Präzision
    ($\mathrm{dps}\approx 30\text{--}50\cdot\lambda$) ist das Bild
    qualitativ anders: $A_\lambda(0.4)\approx 1.004$ ist flach. Die
    methodische Lektion (dps-Bedarf skaliert linear mit $\lambda$ für
    cluster-nahdegenerierte Eigenvektoren) ist selbst substantiell.

  \item \textbf{Roher-$L^2$-Hurwitz-Übergang (falsifiziert 2026-05-15).}
    Der Plan, die Hurwitz-streifenuniforme Konvergenz direkt aus einer
    rohen-$L^2$-Tailabschätzung abzuleiten, scheitert auf zwei Fronten:
    (a) der archimedische Tail hat Norm $O(\lambda^{-1/2})$ --- exakt
    Randexponent $\alpha=0$, keine Marge; (b) der Prim-Dirac-Tail
    divergiert im $m=1$-Branch. Negativ-Sobolev rettet dies nicht. Die
    Diagnostik identifiziert Operatornorm-Topologie (d.\,h.\ Branch~B)
    als einzigen gangbaren Ersatz. Das Connes--Hurwitz-Passage-Projekt
    wurde am 2026-05-17 auf eine rein darstellende Rolle umgewidmet
    (\texttt{Geiger2026hurwitz}, in Vorbereitung).
\end{itemize}

\begin{remark}[Gemeinsames Muster]
\label{rem:spectral-pattern}
Diese sechs Ausschlüsse teilen ein strukturelles Muster: der
Minimaleigenvektor-Cluster von $QW_N$ ist die Quelle der technischen
Schwierigkeit. Methoden, die eine uniforme Lücke, eine
dünn-lokalisierte Eigenfunktion oder ein generisches
Funktionsklassenverhalten voraussetzen oder erfordern, scheitern alle.
Produktive Wege arbeiten \emph{mit} dem Cluster (Waterline-/Cluster-Quotienten-Reformulierungen, wachstumsadaptierte Normalisierungen) statt gegen ihn.
\end{remark}


% =============================================================================
\section{Ausgeschlossene CCM-Wege (Branch~Z-Erbe)}
\label{sec:excluded-ccm}
% =============================================================================

Das CCM-Mikrocluster-Programm hat seinen eigenen Katalog toter Routen
angesammelt. Sie beeinflussen die bedingte Reduktion von Branch~Z nicht.

\begin{itemize}
  \item \textbf{Massebasierte Mikrocluster-Definition.} Die Definition
    des Mikroclusters über ein Massekriterium ist zirkulär gegenüber MS2.
    Ersetzt in Zookeeper v1.2+ durch einen \emph{spektralen} Mikrocluster
    (Niedrigenergie-Eigenraum, begrenzt durch eine Outer Gap $g_*$ der
    Ordnung eins).

  \item \textbf{Naiver Hurwitz-Übergang.} Eine direkte Hurwitz-Anwendung
    ohne Approximationsqualitätskontrolle liefert Aussagen, die für den
    err-collapse-Abschluss zu schwach sind. Ersetzt durch einen
    Slepian--Pollak-/PSWF-basierten Übergang.

  \item \textbf{Anker-Störung / Banach-Kontraktion (C2br$'$,
    falsifiziert).} Eine Banach-Norm-Kontraktion $C_{\max}<M$ scheitert
    um den Faktor 55--7000 auf allen zwölf getesteten Konfigurationen.

  \item \textbf{Projizierte Frontier-Auslöschung $S_{bulk}\approx
    -S_{bd}$ (C2bs, falsifiziert).} Beide Beiträge haben das gleiche
    Vorzeichen und sind beide von der Größenordnung $\sim 10^{-7}$. Der
    wahre Auslöschungsmechanismus liegt in der Eigenwertgleichung
    $(A-w_0)k\approx 0$.

  \item \textbf{Uniformes Massekriterium C2ak.} Die gesamte
    Off-Cluster-Masseschranke divergiert mit $\lambda$
    ($6.68\to 15.3\to 1865$ über $\lambda=3,5,7$), weil sie interne
    Auslöschung verwirft.

  \item \textbf{Kanonische lokale Geometrie C2al (Richtungsdefekt).}
    Die geometrische Ungleichung hat die falsche Richtung --- die
    korrekte Form ist $g_R/d_m\le T_m^\beta/\beta_0$ (untere Schranke
    für $d/g$, nicht obere).

  \item \textbf{Clusterselektion via $\Delta$-Projektion.} Einfache
    $\delta_N$-Projektion ist kein robuster Clusterselektor: bei
    $\lambda=5,N=55$ wählt \texttt{best\_zero=k=16}, während die
    anderen beiden $k=17$ selektieren.

  \item \textbf{Log-Konkavität / detailliertes Gleichgewicht (C2cc
    H1/H2/H4).} Die Hypothesen, dass $\beta_n$ oder $|b_n|$
    log-konkav sind (H1), dass die spektrale Residuenkette eine
    Detailed-Balance-Struktur hat (H2), oder dass ein Mertens-Anker
    gilt (H4), wurden bei
    $\lambda=3,N=30,\mathrm{dps}=50$ getestet und falsifiziert.
    Strukturelle Ursache: arithmetische (primzahlgetriebene)
    Oszillation.

  \item \textbf{Direkte Frontier-Dominance v2.1 (Killer~0, 2026-05-14).}
    Das v2.1-Argument verwendete einen gemeinsamen Rayleigh-Testvektor,
    um $\lambda_{\min}(D_{tot})<0$ zu beweisen, und identifizierte dann
    $\Delta_{low}=\lambda_{\min}(D_{tot})$ mit
    $\Delta=\lambda_1^+-\lambda_1^-$. Das Reviewer-Gegenbeispiel
    $A=\mathrm{diag}(0,10), B=\mathrm{diag}(5,-5)$ zeigt, dass dies
    \emph{nicht} folgt. Even Dominance für asymptotisches $\lambda$
    bleibt bedingt durch eine separate Odd-Sektor untere Schranke.
    (Dieser Punkt ist Even-Sektor-spezifisch und in
    Teil~II~\cite{Geiger2026partII} detailliert dokumentiert; er wird
    hier aufgeführt, weil die Falsifikation auch die vergleichende
    Analyse der Branches~B und~Z umstrukturierte.)

  \item \textbf{Euler-Diagonal-Determinante (Z1, geschlossen durch
    Schatten-Trichotomie, 2026-06-09).} Die Realisierung des
    Zähl-Objekts als regularisierte Determinante über die
    Euler-Diagonale scheitert strukturell: Auf
    $\operatorname{Re} s = 1/2$ liegt die Diagonale in
    $\mathcal{S}_3\setminus\mathcal{S}_2$ ($\det_2$ divergiert via
    Mertens), während $\det_3$ den arithmetischen $k=1,2$-Kern verliert.
    Nur eine andere Singulärwert-Klasse (Ruelle-Transfer-Typ) bleibt
    unerforscht.

  \item \textbf{Kontur-Zählung allein innerhalb von $\Omega_E$ (F3,
    2026-06-09).} Den Nullstellen-Zähler durch Konturen zu schließen,
    die ganz im Bereich etablierter Konvergenz liegen, scheitert an
    einem Konvexitäts-Lemma; jeder Gewinn re-importiert
    Positivitäts-/Tightness-Apriori.

  \item \textbf{Konstantes-$C$ Paket-Rouch\'e / Unendlich-Kamm-Referenz
    (2026-06-10/11).} Auf Paket-Konturen scheitert die
    Rouch\'e-Dominanz gegen ein konstantes Vielfaches von $\Xi$ im
    $N\to\infty$-Limes bei allen getesteten $c$ (Margin-Plateau
    $1{,}541$ bei $c=9$; steigend durch $\approx 0{,}95$ bei $c=14$);
    die Vervollständigung des Kamms zu seiner geschlossenen Sinc-Form
    macht die Referenz strikt schlechter (Typ-Kollision). Die
    $z$-abhängige Normalisierungsfrage ist in \Cref{sec:det-channel}
    geklärt.

  \item \textbf{Wurzel-Paarungs-Bulk-Modelle ($B_K$, PolyB,
    2026-06-09/10).} Den Bulk-Faktor durch Paarung äußerer
    Säkular-Wurzeln gegen äußere Nullstellen zu modellieren ist
    wirkungslos (die Wurzeln treffen die Nullstellen unter Mess-Floor,
    $B_K \equiv 1$); das nackte lokale Wurzelpolynom gegen $C\,\Xi$
    scheitert um zwei Größenordnungen (Margin $35$). Der Bulk wird
    ausschließlich vom Kompensationspolynom und der $\Xi$-Hüllkurve
    getragen.
\end{itemize}

\begin{remark}[CCM-Muster: Auslöschung, nicht Betrag]
\label{rem:ccm-pattern}
Die toten CCM-Routen teilen ein anderes Muster als die spektrale Seite:
jede versucht, eine Größe durch ihren Einzelterm-Betrag zu beschränken,
und jede wird durch große interne Auslöschung besiegt. Die lebenden
Routen faktorisieren die Auslöschung explizit durch
Projektoridentitäten oder skalare säkulare Gleichungen, niemals durch
einfache Normdominanz.
\end{remark}


% =============================================================================
\section{Methodische Sackgassen}
\label{sec:methodical}
% =============================================================================

Eine kleine Anzahl methodischer Routen wurde auf der spektralen Seite
getestet und aus strukturellen Gründen ausgeschlossen.

\begin{itemize}
  \item \textbf{Dobrushin--Zegarlinski-Eindeutigkeit $\to$ RH-CAP.}
    Dobrushin--Zegarlinski-Eindeutigkeit funktioniert nur auf Gittern
    mit lokalisierten Gibbs-Maßen. RH hat weder das eine noch das
    andere; das Fehlen multiplikativer Unabhängigkeit erzwingt eine
    nichtlokale Primzahlkopplungsstruktur, die die
    Dobrushin-Kontraktionssituation zerstört.

  \item \textbf{Sturm-Knotenzählung (naiv).} Klassische
    Sturm--Liouville-Knotenzählung für die Eigenfunktionen von $QW_N$
    scheitert, weil der Operator endlichdimensional ist und die
    gerade/ungerade Spektrallücke subexponentiell klein ist. Die
    Knotenzählung liefert keine Information auf der relevanten
    Präzisionsskala.

  \item \textbf{Naive Abel-Summation für den fernen Tail.} Die
    Anwendung der Abel-Summation auf die fernen Primzahlsummen
    $\sum_{p>\lambda} (\log p)^k / p^s$ ohne Kontrolle der
    arithmetischen Oszillation von $\theta(x)-x$ führt zu divergenten
    Fehlertermen. Das PNT-Transferlemma (Teil~II) zeigt, wie dies
    korrekt durchgeführt wird.

  \item \textbf{Naive Superdeterminantenform für Möbius-$\mathbb{Z}_2$.}
    Die formale Identität
    $\mathrm{sdet}(I - L_s) = \det(I-L_s^+)/\det(I-L_s^-) = 1/\zeta(s)$
    ist algebraisch korrekt für endliche Primzahlmengen, konvergiert aber
    nicht für die volle Primzahlmenge bei $\Re(s) = \tfrac{1}{2}$
    (Spurklassen-Obstruktion K4 aus Teil~I).

  \item \textbf{Friedrichs + Petersson als Hilbert--P\'olya-Hebel.}
    Der Plan, die Friedrichs-Erweiterung des Petersson-Innenprodukts zu
    nutzen, um eine selbstadjungierte Realisierung zu konstruieren,
    deren Spektrum die Riemann-Nullstellen kodiert, scheitert: die
    relevante Bilinearform ist nicht halbbeschränkt auf dem Definitionsbereich, der nötig ist, um die Nullstellen auf der
    kritischen Geraden zu erfassen.
\end{itemize}


% =============================================================================
\section{Update (v3.2): Geschlossene Einheiten, Juni--August 2026}
\label{sec:v32-update}
% =============================================================================

Zwischen Juni und August 2026 wurde eine Reihe endlicher, auditgestützter
Arbeitseinheiten auf und um den spektralen Sektor geschlossen. Sie werden
hier im selben Geist dokumentiert wie die Ausschlusskataloge oben: Jeder
Eintrag nennt, was geprüft wurde, wie weit es verifiziert ist, was es für
die Route bedeutet und --- ausdrücklich --- was er \emph{nicht}
beansprucht. Keine der Einheiten schließt RH, und keine ändert den Status
der bedingten Reduktion von Branch~B und~Z aus den vorangehenden
Abschnitten.

\begin{remark}[Zum Lesen der Zahlenangaben in diesem Abschnitt]
\label{rem:v32-counts}
Zwei Arten von Verhältnisangaben treten unten auf und dürfen nicht gleich
gelesen werden. Eine \emph{Verifikationsquote} --- ,,Produzent $18/18$``,
,,unabhängiger Prüfer $22/22$`` --- nennt Prüfungen, die durchgeführt und
bestanden wurden. Eine \emph{Fehlmengenangabe} nennt dagegen, wie viele
geforderte Posten \emph{fehlen}: Wo unten steht, dass keine der fünfzehn
Voraussetzungen erfüllt ist, dass keine der vierzehn beweistragenden
Kategorien vorliegt oder dass keines der einundzwanzig Pakete existiert,
sind das Angaben über fehlende Posten und ausdrücklich keine
Bestehensquoten. Aus diesem Grund ist jede Fehlmenge ausgeschrieben,
während echte Verifikationsquoten als Zahlen stehen und als solche benannt
sind.
\end{remark}

\subsection{Block A: Die Freedman/Weyl-Kontinuumsroute}
\label{sec:v32-block-a}

\begin{itemize}
  \item \textbf{Quotient$\to$Weyl-Verpflichtungsvertrag V1 (exakter
    endlicher Vertrag mit Gegenmodell).} Die Einheit legt vertraglich fest,
    was ein Lift von der endlichen Quotientenpositivität zur
    Weyl-/Kontinuumsaussage leisten muss, und sichert dies durch ein
    Gegenmodell ab, das den Schluss von Positivität auf $\ker R$ auf volle
    Positivität sperrt. Verifikation: Produzent $14/14$, unabhängig
    $16/16$, wobei ein Byte-Angriff auf die gebundenen Quellen fail-closed
    scheitert. Eine spätere Lektüre des Projektregisters (Stand~V1.1,
    jünger als der Planstand) präzisiert das Bild: Die Quelle erledigt die
    Randpflichten selbst --- $\beta_{\mathrm{bdy}} = 0$ und
    $D_{\mathrm{bdy}} = 0$ per primitiver Endpunktrechnung, und
    $D_q|_Y = 0 \Longleftrightarrow D_q = 0$ auf $X_R$ per Dichte --- und
    stellt die Äquivalenz
    $D_q = 0 \Longleftrightarrow \Gamma^*\Gamma \le C
    \Longleftrightarrow Q_\Phi \ge 0$ als direkten Lift-Satz des Originals
    fest. Für die Routenkarte heißt das, dass der Weyl-Lift auf
    Registerstand~V1.1 keine diffuse Liste von sieben Kontinuumspflichten
    mehr ist, sondern auf eine einzige analytische Restpflicht konzentriert
    ist: $\Gamma^*\Gamma \le C$ auf $X_R$, uniform in $\omega$, samt den
    Rahmenpflichten. Der endliche Quotientenbefund allein trägt den Lift
    weiterhin nicht. Claim-Grenze: Es wird kein PASS für
    $\Gamma^*\Gamma \le C$, für den Lift, für Weyl, für KLM oder für RH
    behauptet, und die ältere Formulierung ,,sieben Kontinuumspflichten
    offen`` ist durch Registerstand~V1.1 überholt.

  \item \textbf{Uniform-$\omega$-Umgewichtungsaudit V1 (endlicher
    Negativbefund).} Positive und lokal uniform integrierbare
    Zweiggewichte erhalten die aggregierte Quadratdomination \emph{nicht}
    ohne Zusatzhypothese. Verifikation: Produzent $15/15$, unabhängig
    $17/17$, fail-closed. Da der Befund endlich und exakt ist statt bloß
    numerische Evidenz, lässt er sich unmittelbar als Routenentscheidung
    verwenden. Seine Folge ist, dass die naheliegende
    Umgewichtungsstrategie für $|\omega| < 1/2$ stirbt; offen bleibt die
    operatorgebundene Alternative (zweigweise Domination, uniform
    gewichteter Intertwiner, konjugierte Kontraktion und Parametercover mit
    positiver Marge). Claim-Grenze: Dies ist \emph{kein} Gegenbeispiel zum
    tatsächlichen Freedman-Kernel --- widerlegt ist die Strategie, nicht
    die Aussage.

  \item \textbf{Endpunkt$\to$Kontinuums-Folgekettenvalidierung und
    DAG-Audit (negativ).} Elf Dashboardzeilen wurden auf vier
    Quellenkohärenz-Schnitte reduziert, und die Folgekette Endpunkt $\to$
    Rang $\to$ Green $\to$ Unique Continuation $\to$ Active Range wurde
    negativ validiert. Die gebundenen Quellen führen für diese Einheit
    keine Mutationszahlen, und es werden hier auch keine ergänzt. Der
    Befund ist kein universelles No-Go: Er gilt negativ für genau diese
    Kette und für keinen weiteren Anspruch. Als Nebenergebnis wurde der
    Waisenakten-Querbezug als endliches Diagnoseinstrument geschlossen.
    Claim-Grenze: Weder die Endpunktschicht noch das Kontinuumsziel ist
    widerlegt; was scheitert, ist die spezifische Verkettung beider.

  \item \textbf{Ganzzeit-/Kontinuums- und Pullback-Readiness-Audit V1
    (fail-closed).} Fünfzehn Voraussetzungen einer ganzzeitfähigen,
    kontinuumsfähigen und nichtzirkulären
    KLM $\to$ de-Branges-Route wurden auditiert, und \emph{keine der
    fünfzehn ist erfüllt} --- eine Fehlmengenangabe im Sinne von
    \Cref{rem:v32-counts}, keine Bestehensquote. Das Audit selbst ist
    sauber und führt eigene Quoten: Produzent $18/18$, unabhängiger Prüfer
    $22/22$, Terminalstatus \texttt{NOT\_\allowbreak READY\_\allowbreak
    MISSING\_\allowbreak BOUND\_\allowbreak CONTINUUM\_\allowbreak
    AND\_\allowbreak PULLBACK\_\allowbreak OBJECTS\_\allowbreak
    NEW\_\allowbreak MATHEMATICAL\_\allowbreak INPUT\_\allowbreak
    REQUIRED}. In diesem Pfad verbleibt kein autonom ausführbarer
    statischer oder endlicher Rechenschritt, sodass Fortschritt neuen
    mathematischen Input erfordert statt weiterer Rechnung. Claim-Grenze:
    Die Brücke ist unfertig, nicht widerlegt --- das Audit misst
    Bereitschaft, nicht Wahrheit.
\end{itemize}

\subsection{Block B: Die KLM/de-Branges-Route}
\label{sec:v32-block-b}

\begin{itemize}
  \item \textbf{KLM $\to$ de-Branges-Pull\-back\-pflicht-Audit V1
    (Kreuz\-term-Gegen\-modell).} Aus separat nichtnegativen Tail- und
    Boundary-Differenzen folgt \emph{nicht} volle signierte
    Gram-Positivität. Verifikation: Produzent $17/17$, unabhängig $19/19$,
    fail-closed. Die Routenfolge ist, dass die Strategie ,,separat
    nichtnegativ genügt`` stirbt und die unbehandelten Kreuzterme damit als
    Ent\-schei\-dungs\-ort für jeden Nach\-folge\-versuch benannt sind.
    Claim-Grenze: Der tatsächliche Shifted-$\Xi$-Kern ist nicht widerlegt;
    es fällt das Hin\-läng\-lich\-keits\-argument, und es fällt an einem
    expliziten Gegenmodell statt an einer numerischen Marge.

  \item \textbf{KLM-Träger-, Isometrie- und Defektlift-Schicht (vier
    geschlossene Befunde).} Positiv: Ein gewichteter Vollgeraden-Träger
    $H_\rho$ realisiert beide Mellin-Hälften (Produzent $28/28$,
    unabhängig $34/34$), und ein zielunabhängiger Fixed-Trace-Träger
    $\mathrm{Mu}$ wurde konstruiert (Produzent $19/19$; Prüfer mit $49$
    Direktprüfungen plus $21/21$). Negativ: Kein Skalar $a_t T_t$ ist
    isometrisch für $t \neq 0$, und ein exakter Gewichtskokzyklus zerstört
    die $k_z$-Eigenrelation (Produzent $22/22$, unabhängig $28/28$); eine
    universelle gemeinsame Kreuz-Gram-Reparatur der Defektlifts
    $J_{(\sigma,u)}$ ist widerlegt (Produzent $26/26$, unabhängig $31/31$,
    mit $126+378$ Kontrollen); und die
    $(\Lambda_a, \mathrm{Mu}_z)$-Nullraumannihilation ist ab Sättigung
    vakuos (Produzent $17/17$, unabhängig $21/21$). Die Routenlesart ist,
    dass die Träger existieren, während die Isometrieroute und die
    universelle Defektlift-Reparaturroute geschlossen sind --- was zugleich
    einen früheren, als Fortschritt verbuchten Posten entwertet.
    Claim-Grenze: Jede Negativaussage ist an ihre genannte Klasse gebunden
    (skalare Isometrien, exakte Kokzyklen, universelle gemeinsame
    Reparaturen), und keine von ihnen widerlegt das KLM-Pullbackziel
    selbst.

  \item \textbf{Vollständiger normalisierter KLM-Residualblock ---
    gerichtet nicht-PSD (Endbefund).} Fünf präregistrierte Hardy-Knoten und
    alle zehn Riesz-Vektoren stammen aus gerichteten normalisierten
    $79\times 10$-Zweigabbildungen, und der Hardy-Zielblock wurde mit
    Arb/Acb neu eingeschlossen ($20$ Moden, Majoranten unter $2^{-120}$).
    Für den Residualblock $\Delta = H - C$ haben alle zehn reellen
    Diagonalboxen strikt negative Oberränder, der größte davon
    $\approx -0.2665171987\ldots$. Verifikation: $26/26$ Quellen
    bytegebunden, $384$-Bit-Produzent $17/17$, unabhängiger
    $512$-Bit-Prüfer $20/20$, fail-closed. Diese gemeinsame
    zweigsymmetrische fünfknotige Polar- und Graphnormierung stirbt damit
    rigoros als positiver Pullback-Kandidat, und zwar als Endbefund statt
    als vorläufiger; zwei gröbere Varianten waren zuvor bereits numerisch
    ausgeschlossen worden (Plus-Matching $(5,0,5)$ mit
    $\lambda_{\min} \approx -0.11543$, Einheits-Direktsummen $(4,1,5)$ mit
    $\lambda_{\min} = -12.2767\ldots$). Claim-Grenze: Das Resultat bindet
    diesen einen endlichen Kandidaten; der Mellin-Split ist dadurch nicht
    widerlegt.

  \item \textbf{Operatorsemantischer Alternativen-Audit V1.} Drei
    zugelassene Klassen wurden auf eine target-freie Alternative geprüft:
    zweiggewichtete Graphmetriken, feste Trace-Augmentationen und
    gekoppelte Defektkanäle. Keine der drei liefert eine eindeutige
    target-freie Wahl, und es wurde keine Matrix erzeugt. Verifikation:
    Produzent $15/15$, Prüfer $17/17$, Terminalstatus
    \texttt{NO\_\allowbreak TARGET\_\allowbreak FREE\_\allowbreak
    UNIQUE\_\allowbreak FINITE\_\allowbreak ALTERNATIVE\_\allowbreak
    MATERIALIZED\_\allowbreak SWITCH\_\allowbreak TO\_\allowbreak
    CONTINUUM\_\allowbreak PULLBACK}. Die Routenfolge ist, dass der finite
    Matrixpfad erschöpft ist, was diese Linie in das Readiness-Audit von
    Block~A oben überführt. Claim-Grenze: Erschöpfung wird allein für die
    drei zugelassenen Klassen behauptet --- eine Klasse außerhalb dieser
    Zulassung ist ungeprüft, nicht ausgeschlossen.
\end{itemize}

\subsection{Block C: Weil/CvS-Numerik}
\label{sec:v32-block-c}

\begin{itemize}
  \item \textbf{RH-V4R-Fremdoperatorlauf (terminal, unabhängig
    verifiziert).} Ein autorisierter Lauf eines Fremdoperators gab
    Rückgabecode~$0$ zurück, mit Checkpoints bei $N = 48/56/64$. Sein
    Terminalergebnis ist \texttt{C13\_\allowbreak FIXED\_\allowbreak
    LADDER\_\allowbreak COMPLETE\_\allowbreak TAIL\_\allowbreak
    GATES\_\allowbreak NOT\_\allowbreak REACHED}: Der Rayleigh-Tail
    besteht, der Fit-Tail scheitert ($0.1273\ldots > 0.0005$). Die
    Verifikation war auf zwei Ebenen unabhängig --- der Ergebnisprüfer
    verwarf $10/10$ Mutationen, der Terminalumschlag-Prüfer $17/17$
    einschließlich einer Checkpoint-Bytemutation. Die Routenlesart ist,
    dass diese Linie gemessen, nicht bestätigt wurde; der offene Punkt
    \texttt{CVS\_\allowbreak N\_\allowbreak RESOLUTION\_\allowbreak
    UNRESOLVED} steht unverändert. Claim-Grenze: Der gescheiterte Fit-Tail
    ist ein Gate-Ergebnis und keine Widerlegung der zugrunde liegenden
    Route, und ohne neuen Vertrag ist kein Retry zulässig.

  \item \textbf{K1-Skalaraudit, transitiver Quellen-Audit und
    fail-closed Vollreplay-Eingabevertrag.} K1 bleibt auf der
    Ursachenebene offen. Das Skalaraudit ($256$-Bit-Arb) ergibt einen
    Fitfehlerquotienten von $3.9533\ldots$, während die Rayleigh-Seite nur
    $1.4489\times 10^{-8}$ spannt. Der transitive Quellen-Audit zeigt, dass
    der Wrapper den V4-Kern zwar ausführt, aber eine Provenienzlücke lässt,
    weil Pfad und SHA in den Manifesten fehlen --- ein historischer Mangel,
    der den Lauf nicht automatisch invalidiert. Ein fail-closed
    Vollreplay-Eingabevertrag benennt nun exakt vierzehn beweistragende
    Kategorien für $N = 40,\dots,64$, und \emph{keine dieser vierzehn
    Kategorien liegt derzeit vor}; im Sinne von \Cref{rem:v32-counts} sind
    die vierzehn die fehlenden Posten, keine bestandenen Prüfungen.
    Claim-Grenze: Die Einheit beziffert den Preis einer Wiederaufnahme von
    K1 exakt, nimmt sie aber weder wieder auf noch hebt sie einen Claim an.

  \item \textbf{Matrix-Produktionslauf D-20260813-003 (terminal
    verifiziert, Kandidat fehlt).} Ein einmalig autorisierter
    Arb-Matrixlauf ($c = 13$, $N = 48/56/64$, $T = 200$, $3225.75$~s) gab
    Rückgabecode~$0$ zurück und bestand die Verifikation. Der Prüfer
    rekonstruierte die Matrizen $49\times 49$, $57\times 57$ und
    $65\times 65$, bestätigte $10/10$ Primkomponenten und verwarf $12/12$
    Mutationen; das Paket ist byteidentisch (SHA
    \texttt{78ba72b5\ldots aab0e7}). Der Terminalstatus lautet
    \texttt{RIGOROUS\_\allowbreak MATRIX\_\allowbreak
    PRIMITIVES\_\allowbreak READY\_\allowbreak CANDIDATE\_\allowbreak
    MISSING}, festgehalten mit \texttt{candidate\_\allowbreak
    verified\_\allowbreak here=false} und \texttt{claim\_\allowbreak
    upgrade=false}. Was der Lauf liefert, sind rigorose Matrixprimitive;
    was er nicht liefert, ist der Kandidat --- Kandidatenrealisierung,
    Winkelkontrolle und die Sektorbedingung
    $\mu_{1,g^+} < \min(\mu_{2,g^-},\, \mu_{1,u^-})$ sind eigene Gates, die
    hier nicht ausgeführt wurden. Claim-Grenze: Primitive bereit ist nicht
    Kandidat verifiziert, und der Lauf hebt ausdrücklich keinen Claim an.
\end{itemize}

\subsection{Block D: Diagnostik und Hygiene}
\label{sec:v32-block-d}

\begin{itemize}
  \item \textbf{CCM--Groskin-Metadatenpaket.} Einundzwanzig normalisierte
    Zellen des Metadatenpakets wurden auditiert. Die gebundenen Quellen
    führen für diese Einheit keine Produzenten- oder Prüferquoten, und es
    werden hier auch keine ergänzt. Das Paket ist ausdrücklich
    \texttt{DIAGNOSTIC\_\allowbreak ONLY\_\allowbreak NOT\_\allowbreak
    VALIDATED}: Es trägt diagnostischen Wert, aber keine Evidenzstufe, und
    seine Negativblöcke sind Warnsignale statt nachgewiesener
    Spektralverschmutzung. Offen sind unter anderem die
    $c = 100$-Runner-Bindung, die Rohspektren, die kleinste-positive
    Zweigauswahl, die Klassifikation, der Spektralverschmutzungstest, die
    Operatorrealisierung, die Grenztopologie und der
    $c/N/T$-Übergang. Claim-Grenze: Nichts in diesem Paket stützt oder
    widerlegt eine spektrale Aussage, und seine Negativblöcke als erwiesene
    Verschmutzung zu lesen hieße, das Material zu überinterpretieren.

  \item \textbf{Residuen-/Determinanten-Hygienegate
    (Drei-Achsen-Labelaudit).} Ein Drei-Achsen-Labelaudit über $13/13$
    Claim-Zeilen, $7/7$ aktive Akten und $21/21$ Achsen, mit $14/14$
    verworfenen unabhängigen Mutationen, etablierte eine Buchführungsregel
    mit Kill-Wirkung. Die Regel hält drei Kanäle strikt getrennt: das
    numerische Eigenresiduum, das meromorphe Nullstellenresiduum und den
    Determinantenkanal aus \Cref{sec:det-channel}. Neue Spektralakten
    müssen diesen Vertrag fortschreiben, statt ihn nur zu wiederholen.
    Seine operative Folge ist die Unterscheidung, dass ein technischer PASS
    kein analytischer Kanal-PASS ist --- ein Label, das die drei vermengt,
    ist unter dem Gate ungültig, unabhängig von den Zahlen dahinter.
    Claim-Grenze: Das Gate ist ein Hygieneinstrument, es entscheidet also
    über die Zulässigkeit von Labels und nicht über die Wahrheit einer
    spektralen Aussage.
\end{itemize}

\subsection{Aktive Fronten (nur Statuszeilen)}
\label{sec:v32-fronts}

Die folgenden Fronten sind zum Zeitpunkt dieses Updates offen. Jede erhält
eine Statuszeile und keine weiteren Details; die Substanz tragen die
Arbeitsregister.

\begin{itemize}
  \item \textbf{Kongruenz-Reduktion / Restgate G-A/G-B.} Der Struktursatz
    liegt vor und das Gate ist definiert, aber nicht ausgeführt, und es ist
    ein neuer Vertrag mit gerichteter Arithmetik nötig.
  \item \textbf{21-Knoten-Produktionslauf.} Keines der einundzwanzig
    Pakete existiert, und ein eigener Produktionsvertrag samt Autorisierung
    steht aus (der Entwurf verwendet $21$ Kollokationsknoten bei Ordnung
    $200$, nicht $200$ Knoten).
  \item \textbf{N1-/$\Phi_T$-Front.} Eine separate Front mit eigenem
    Register; hier nicht zusammengefasst.
  \item \textbf{KLM-Zweigabbildungen.} Bis zum normalisierten
    Residualbefund oben fortgeschritten, ohne autonom ausführbaren Schritt
    ohne neuen mathematischen Konstruktionsinput.
\end{itemize}


% =============================================================================
\section{Zusammenfassung}
\label{sec:summary}
% =============================================================================

\begin{center}
\renewcommand{\arraystretch}{1.3}
\resizebox{\textwidth}{!}{%
\begin{tabular}{l|c|l}
\textbf{Beitrag} & \textbf{Status} & \textbf{Ort} \\\hline
I-1 Normal-Family Reframe: (ii) $\to$ (ii-a) + (ii-c) & rigoros & \Cref{sec:i1-reframe} \\
(ii-c) absorbiert in (ii-a) unter (Z) & rigoros & \Cref{sec:i1-reframe} \\
Foundation-Lock: kanonisches Kriteriumsobjekt & rigoros & \Cref{sec:i1-reframe} \\
Konvergierte Empirie: $A_\lambda(0.4) \approx 1.004$ (7 Punkte) & numerische Evidenz & \Cref{tab:route-b-empirics} \\
Multiplikative kanonische Form & empirisch (0.012\% Anpassung) & \Cref{sec:canonical-form} \\
Theorem~B-1: $A_\lambda \leq 1.0063$ (bedingt) & bedingt & \Cref{thm:B1} \\
$c_\Xi$ Hadamard-Krümmung (RH-neutral) & analytisch & \Cref{sec:spectral-constants} \\
CCM: drei C2ca-Eingaben identifiziert & bedingte Reduktion & \Cref{sec:c2ca-inputs} \\
CCM: effektive Mikroclusterlücke $g_{\text{eff}} \approx 5\text{--}6$ & numerische Evidenz & \Cref{sec:branch-z-empirics} \\
Determinanten-Kanal: exakte Buchhaltung, nackte $P\!\leftrightarrow\!\Xi$-Paarung, Rate $\sigma_1(\gamma>T_{\mathrm{hit}})$ & numerische Evidenz & \Cref{sec:det-channel} \\
6 ausgeschlossene spektrale Wege + 13 ausgeschlossene CCM-Wege & dokumentiert & \Cref{sec:excluded-spectral,sec:excluded-ccm} \\
5 methodische Sackgassen & dokumentiert & \Cref{sec:methodical} \\
\end{tabular}%
}
\end{center}

\noindent Der spektrale Sektor liefert zwei unabhängige bedingte
Reduktionen der RH, jede mit einer gut charakterisierten offenen
Fläche. Branch~B hat die kleinste offene Fläche (eine einzige
Hypothese~(ii-a)) und die stärkste empirische Unterstützung
($A_\lambda \leq 1.0063$ bedingt durch uniforme
Log-Krümmungskontrolle). Die routenübergreifende Synthese aller drei
Branches wird in Teil~IV~\cite{Geiger2026partIV} entwickelt; der
kondensierte Atlas und Ausblick in Teil~V~\cite{Geiger2026partV}.


\bigskip
\noindent\textbf{Danksagungen.} Berechnungen wurden auf einem
Mac~Studio (Apple M1~Max, 32~GB) durchgeführt. Die
Intervallarithmetik-Bibliothek mpmath wurde von Fredrik Johansson
entwickelt.


\section*{KI-Offenlegung}

Dieses Manuskript wurde in umfassender Zusammenarbeit mit
KI-Sprachmodellen (Claude, Anthropic; Copilot, Microsoft/GitHub; Gemini,
Google DeepMind) entwickelt. Die KI trug zur konzeptionellen Erkundung,
analytischen Arbeit, Literaturrecherche, zum Schreiben und zur
kritischen Überprüfung bei. Der Autor, Lukas Geiger, konzipierte die
Forschungsrichtung, brachte Domänenwissen ein und übte die endgültige
redaktionelle Verantwortung aus. Der Autor übernimmt die volle und
alleinige wissenschaftliche Verantwortung für alle Inhalte,
einschließlich etwaiger Fehler.


\bibliographystyle{plain}
\bibliography{fst-rh-references}

\end{document}
